\section{Introduction}
\label{sec:tpc25-introduction}

The primitive M\"obius-tail program of TPC-18--24 reorganizes a
fixed-shift prime-pair residual into dyadic row packets
\(m=\ell d\), with a large prime source \(\ell\) and a shorter
squarefree factor \(d\) \citep{WangTPC18,WangTPC19,WangTPC20,WangTPC21,WangTPC22,WangTPC23,WangTPC24}.
TPC-24 reaches a genuine one-sided divisor shell
\[
 u\le S<v\le T
\]
in which the new endpoint \(T\) may lie beyond the Poisson length
\(J\).  Its nonzero additive frequencies are controlled by an
asymmetric exact-conductor rectangle.  The physical additive zero
frequency is treated separately through the shared divisor
\(g=(u,v)\mid m-n\).  That argument gives, uniformly on every retained
off-diagonal row pair,
\begin{equation}
 K^0_{S,T}(m,n)
 \ll_{A,\kappa}
 (\log X)^{-A}+X^{-s+\kappa+o(1)},
 \qquad S=X^{s+o(1)},\quad 0<\kappa<s.
 \label{eq:tpc25-inherited-pointwise}
\end{equation}
The first term comes from cancellation between the two endpoints of a
complete M\"obius shell; the second comes from large shared factors of
the nonzero row determinant.

The remaining issue is not hidden in \eqref{eq:tpc25-inherited-pointwise}.
TPC-24 states its separated row interface with
\[
 \|x\|_2^2+\|y\|_2^2\ll QX^\eps,
 \qquad
 \sum_{m,n}|x_my_n\mathfrak m(m,n)|\ll Q^2X^\eps.
\]
Multiplying the first term of \eqref{eq:tpc25-inherited-pointwise} by
\(X^\eps\) does not produce decay for fixed \(A\) and \(\eps\).
Accordingly TPC-24 keeps a separate fixed-polylogarithmic projective-
mass hypothesis, or else assumes that the test function has zero
physical mean.  Its conclusion section isolates two possible next
steps: obtain a row-averaged operator estimate, or recover a
fixed-polylogarithmic projective decomposition from the actual packet.

The main point of the present paper is that the physical packet should
not be separated before frequency zero is extracted.  Before that
separation, the opened-\(d\) quadratic form has the explicit one-row
provenance
\begin{equation}
 w_{\ell,d}=\mu(d)(\log\ell)\times
 \text{a bounded dyadic factor}.
 \label{eq:tpc25-row-provenance-preview}
\end{equation}
For fixed rows, Poisson frequency zero merely integrates the uniformly
bounded orbit weight.  The resulting two-row scalar is still bounded,
and Chebyshev's estimate gives
\begin{equation}
 \sum_{\ell,d}|w_{\ell,d}|\ll LD=Q.
 \label{eq:tpc25-row-mass-preview}
\end{equation}
Thus the physical zero-frequency pair mass is \(O(Q^2)\).  This is a
provenance statement about the actual TPC-18 packet, not a consequence
of the abstract soft \(L^2\) interface.

\subsection{Main results}

The paper has three logically distinct conclusions.

First, the zero-before-separation theorem inserts
\eqref{eq:tpc25-row-mass-preview} directly into
\eqref{eq:tpc25-inherited-pointwise}.  For one actual opened-\(d\)
packet it gives
\begin{equation}
 |\cZ^0_{S,T}|
 \ll JQ^2\left\{(\log X)^{-A}
       +X^{-s+\kappa+o(1)}\right\}.
 \label{eq:tpc25-zero-preview}
\end{equation}
Fixed bounded-overlap dyadic, content, and residue reassembly costs at
most a fixed logarithmic power.  The calibration exponent \(A\) is
then chosen after that cost.  Consequently the extra projective-mass
hypothesis of TPC-24 is verified, rather than assumed, for the actual
one-sided packet.

Second, combine \eqref{eq:tpc25-zero-preview} with the nonzero-
conductor theorem of TPC-24.  If
\[
 Q=X^{\beta+o(1)},\quad J=X^{1-\beta+o(1)},
 \quad S=X^{s+o(1)},\quad T=X^{t+o(1)}
\]
and
\begin{equation}
 \eta_{\rm asym}(\beta;s,t)
 =\min\left\{
 1-\frac{3\beta}{2},
 \frac{\beta+1-2s-2t}{2},
 \frac{3-\beta-s-5t}{4}
 \right\}>0
 \label{eq:tpc25-eta-preview}
\end{equation}
with fixed margin, then the actual calibrated mixed increment has
arbitrary logarithmic saving under the remaining inherited interfaces.
This conclusion allows a general nonzero-mean fixed smooth test within
the factorable provenance class.  At \(S=R=X^{1/2-\delta+o(1)}\) and
\(T=RX^{\xi+o(1)}\), it retains the strict window \(R<J<T\) already
identified by TPC-24.

Third, we prove that this positive result cannot be promoted to
arbitrary arrays satisfying only the soft row interfaces.  The shared
factor \(g=1\) has no row oscillation.  A constant off-diagonal kernel
whose entries are \(\exp(-\sqrt{\log X})\) is smaller than every fixed
negative power of \(\log X\), yet suitable \(X^{o(1)}\)-sized row
weights make its bilinear form saturate \(Q^2\).  This abstract
countermodel is reflected by a finite exact M\"obius kernel:
\[
 H=1,\quad R=S=2,\quad T=3
 \quad\Longrightarrow\quad
 K^0_{2,3}(m,n)=\frac{(\log2-2)\log3}{6}
\]
for every pair of distinct prime rows larger than \(3\).  Hence the
operator norm and Hilbert--Schmidt norm are both of order the number of
rows; a direct Hilbert--Schmidt argument has no dimensional gain.

We finally identify the complete low-row-frequency geometry in a
prime-row no-long-wrap model.  The \(g\)-layer is constant on residue
classes modulo \(g\), so the correlated part of the kernel lies in the
span of residue-class indicators.  Projecting those spaces away gives
an exact deflated operator estimate.  This is a concrete route for a
future both-new-shell problem, but the generic Schur mask and the
general \(d\)-block geometry must then be retained explicitly.

\subsection{Logical scope}

The positive closure in this paper concerns the physical provenance
class \eqref{eq:tpc25-row-provenance-preview}.  It does not apply to an
arbitrary uniformly smooth family \(F_{\alpha_1,\alpha_2}\) with
arbitrary coefficients \(z_{\alpha_1,\alpha_2}\), nor does it prove
that every energy-bounded row vector is balanced against the principal
row mode.  Conversely, the energy-only countermodel is a logical
non-implication, not a lower bound for the wide M\"obius shells in the
target exponent region.

Only one mixed component
\(P_{m,S}(P_{n,T}-P_{n,S})\), and its transpose, is closed here.  The
both-new square, the ultra-long complement fibers, the full residual
reassembly, positivity, and a prime-pair main term remain separate
stages.  In particular no assertion below is a Hardy--Littlewood
asymptotic, a twin-prime theorem, or a general violation of the sieve
parity barrier.
